\documentclass[10pt,a4paper]{amsart}
\usepackage[margin=19mm]{geometry}
\usepackage{amsmath,amssymb,mathtools}
\usepackage[scr=rsfs,cal=euler]{mathalfa}
\usepackage{xfrac}
\usepackage[unicode,hidelinks,bookmarks=false]{hyperref}
\newcommand{\ie}{i.e.\ }
\newcommand{\cf}{cf.\ }
\newcommand{\resp}{respectively\ }
\newcommand{\Subsection}{\subsection}
\providecommand{\nobreakdash}{\nobreak\hbox{-}}
\setlength{\emergencystretch}{2em}
\multlinegap0pt
\usepackage{amssymb}
\usepackage{float}
\newtheorem{theorem}{Theorem}
\title{Eigenvalues of Hypergraph Products and Reciprocal Eigenvalue Property}
\author{Shashwath S Shetty, K Arathi Bhat}
\date{Source: arXiv:2604.23365v1 (CC BY 4.0)}
\begin{document}
\maketitle

\noindent\emph{Source and licence.} Eigenvalues of Hypergraph Products and Reciprocal Eigenvalue Property. Version arXiv:2604.23365v1 (CC BY 4.0), distributed by arXiv under the Creative Commons Attribution 4.0 International licence, http://creativecommons.org/licenses/by/4.0/, as recorded in the arXiv licence metadata for this exact version and shown on the abstract page https://arxiv.org/abs/2604.23365v1. Full licence notice: \texttt{licenses/arxiv-2604.23365-CC-BY-4.0.txt}.\par\smallskip

\noindent\textit{Editorial scope.} Counted unit: the article's theorem theorem (source0.tex lines 554--564). The article's complete original proof of it is delivered with it (source0.tex lines 565--597, one \texttt{proof} environment of 33 lines). No mathematics is written, repaired or re-derived: the statement and the proof are the article's own lines, in the article's own notation, and no retained line is substituted or re-flowed.  Retained uncounted as the source-local context the proof needs: lines 263--265.  Nothing was omitted from any retained span, so the package carries no omission record.  No retained line contains a citation, so the wrapper carries no bibliography and no bibliography entry.  No retained line contains a figure environment, an image-inclusion command or drawing code, so no image is delivered and the package declares no asset dependency.  Editorial formatting only, in addition: an AMS \texttt{amsart} preamble that restates the article's own declarations and macros used by the retained lines, namely the environments \texttt{theorem} on separate counters, together with the standard packages those lines need.  The retained environments print in this document as Theorem 1 in order of appearance, whereas the article prints the same items with the numbering of its own full document.  The complete native source of the article as distributed in the arXiv e-print for this exact version is bundled unchanged as \texttt{sources/199157/source0.tex}.
		\begin{equation}\label{non-main_condition}
			\sum\limits_{S \subset [n], |S|=r-1} \mathbf{x}^S = 0.
		\end{equation}

	\begin{theorem}
		Given an $r$-uniform hypergraph $\mathcal{H}_1$ on $n_1$ vertices and a $d$-regular, $r$-uniform hypergraph $\mathcal{H}_2$ on $n_2$ vertices, let $\mathcal{H}= \mathcal{H}_1 \odot \mathcal{H}_2. $
		\begin{itemize}
			\item If $\lambda$ is an eigenvalue of $\mathcal{A}(\mathcal{H}_1)$, then $\mu_i ~(\neq d)$ are the eigenvalues of $\mathcal{A}(\mathcal{H})$, 
			where $\mu_i, 1\leq i \leq r$ are the roots of the polynomial $(\mu -\lambda)(\mu-d)^{r-1} - \binom{n_2}{r-1}\binom{n_2-1}{r-2}^{r-1}$.
			Furthermore, if $\eta_i$ is a non-main eigenvalue of $\mathcal{A}(\mathcal{H}_2)$, then $\eta_i$ is also an eigenvalue of $\mathcal{A}(\mathcal{H})$ with Hspan multiplicity at least $n_1$.
			\item If $\beta$ is an eigenvalue of $\mathcal{Q}(\mathcal{H}_1)$, then $\mu_i ~(\neq 2d+\binom{n_2-1}{r-1})$ are the eigenvalues of $\mathcal{Q}(\mathcal{H})$, where $\mu_i, 1\leq i \leq r$ are the roots of the polynomial 
			$$\left(\mu-\beta-\binom{n_2}{r-1}\right)\left(\mu-2d-\binom{n_2-1}{r-1}\right)^{r-1}-\binom{n_2}{r-1}\binom{n_2-1}{r-2}^{r-1}.$$
			Furthermore, if $\alpha$ is a non-main eigenvalue of $\mathcal{Q}(\mathcal{H}_2)$, then $\alpha + \binom{n_2-1}{r-2}$ is also an eigenvalue of $\mathcal{Q}(\mathcal{H})$ with Hspan multiplicity at least $n_1$.
		\end{itemize}
	\end{theorem}

	\begin{proof}
		Let $[(n_2+1)n_1]$ be the vertex set of $\mathcal{H} = \mathcal{H}_1 \odot \mathcal{H}_2$, the corona of $\mathcal{H}_1$ with $\mathcal{H}_2$ ($d$-regular). Without loss of generality, assume that $[n_1]$ is the vertex subset corresponding to $\mathcal{H}_1$ and $\{  n_1 + (j-1)n_2+ 1, n_1 + (j-1)n_2+ 2,\ldots, n_1+jn_2\}$ be the vertices corresponding to the $j^{th}$ ($1 \leq j \leq n_1$) copy of $\mathcal{H}_2$ in $\mathcal{H}$. Suppose that $(\lambda,\mathbf{x})$ is an eigenpair of $\mathcal{A}(\mathcal{H}_1)$, and $\mu=\mu_i \neq d$ are the roots of the polynomial $(\mu-\lambda)(\mu-d)^{r-1}-\binom{n_2}{r-1}\binom{n_2-1}{r-2}^{r-1}$.
		Define $\mathbf{y}=\mathbf{y}^{(i)}:=\left[\mathbf{x}^\intercal, \frac{\binom{n_2-1}{r-2}x_1}{\mu_i-d} \mathbf{1}_{n_2}^\intercal, \frac{\binom{n_2-1}{r-2}x_2}{\mu_i-d} \mathbf{1}_{n_2}^\intercal, \ldots,  \frac{\binom{n_2-1}{r-2}x_{n_1}}{\mu_i-d} \mathbf{1}_{n_2}^\intercal \right]^{\intercal}$.
		We claim that $(\mu, \mathbf{y})$ 
		%(for the sake of simplicity we use $(\mu, \mathbf{y})$  instead of ($\mu_i, \mathbf{y}^{(i)}$)) 
		forms an eigenpair of $\mathcal{A}(\mathcal{H})$. 
		For any $j \in [n_1]$, we have 
		\begin{align*}
			\mu_i y_j^{r-1}- (\mathcal{A}(\mathcal{H})\mathbf{y}^{r-1})_j 
			%&=  \mu y_j^{r-1}-\sum\limits_{e \in \mathcal{E}(\mathcal{H}_1)} \mathbf{y}^{e\setminus \{ j \}} - \sum\limits_{e \in \mathcal{E}(\mathcal{H}_1) \setminus \mathcal{E}(\mathcal{H}_1)} \mathbf{y}^{e\setminus \{ j \}} \\
			&=  \mu_i x_j^{r-1}-\sum\limits_{e \in \mathcal{E}(\mathcal{H}_1)} \mathbf{x}^{e\setminus \{ j \}} - \sum\limits_{\substack{S \subseteq \{ n_1 + (j-1)n_2+ 1,\ldots, n_1+jn_2\}\\ |S|=r-1}} \mathbf{y}^{S} \\
			&=  \mu_i x_j^{r-1}-\lambda x_j^{r-1} - \frac{\binom{n_2}{r-1}\binom{n_2-1}{r-2}^{r-1} x_j^{r-1}}{(\mu_i-d)^{r-1}} =0.
		\end{align*}
		For any $k \in \{ j n_1+ 1, j n_1+2,\ldots, (j+1)n_1\},$ $j \in [n_1]$, we have 
		\begin{align*}
			\mu_i y_k^{r-1} - (\mathcal{A}(\mathcal{H})\mathbf{y}^{r-1})_k 
			&= \mu_i \frac{\binom{n_2-1}{r-2}^{r-1} x_j^{r-1}}{(\mu_i-d)^{r-1}} - \sum\limits_{e \in \mathcal{E}(\mathcal{H}_1)} \mathbf{y}^{e\setminus \{ j \}} - \sum\limits_{e \in \mathcal{E}(\mathcal{H}) \setminus \mathcal{E}(\mathcal{H}_1)} \mathbf{y}^{e\setminus \{ j \}} \\
			% &= \mu_i \frac{\binom{n_2-1}{r-2}^{r-1} x_j^{r-1}}{(\mu-d)^{r-1}} - \sum\limits_{e \in \mathcal{E}(\mathcal{H}_2^{(j)})} \mathbf{y}^{e\setminus \{ j \}}- \sum\limits_{\substack{S \subseteq \{ n_1 + (j-1)n_2+ 1, n_1 + (j-1)n_2+ 2,\ldots, n_1+jn_2 \} \setminus \{k \}\\ |S|=r-2}} \mathbf{y}^{S} x_j \\
			&= \mu \frac{\binom{n_2-1}{r-2}^{r-1} x_j^{r-1}}{(\mu-d)^{r-1}}- \frac{d\binom{n_2-1}{r-2}^{r-1} x_j^{r-1}}{(\mu-d)^{r-1}} - \frac{\binom{n_2-1}{r-2}^{r-1} x_j^{r-1}}{(\mu-d)^{r-2}} = 0.
		\end{align*}
		Furthermore, if there exist an eigenvector $\mathbf{z}$ of $\mathcal{A}(\mathcal{H}_2)$ associated with the eigenvalue $\eta$ satisfying Equation \ref{non-main_condition}, then we can see that $\eta$ is also an eigenvalue of $\mathcal{A}(\mathcal{H})$ with corresponding eigenvectors $\left[\begin{matrix}
			\mathbf{0} \\ \mathbf{z} \otimes \mathbf{e_j} 
		\end{matrix}\right], 1 \leq j \leq n_1$, where $e_j, 1 \leq j \leq n_1$ are the canonical vectors of length $n_1$.  
		
		Now, suppose that $(\beta, \mathbf{w})$ is an eigenpair of $\mathcal{Q}(\mathcal{H}_1)$, and $\mu = \mu_i~ ( \neq 2d+\binom{n_2-1}{r-2} )$ is a root of the polynomial \\ $\left( \mu-\beta-\binom{n_2}{r-1} \right)\left( \mu-2d-\binom{n_2-1}{r-2} \right)^{r-1}-\binom{n_2}{r-1} \binom{n_2-1}{r-2}$. Define $\mathbf{y}=\mathbf{y}^{(i)} := \left[ \mathbf{w}^\intercal, \frac{\binom{n_2-1}{r-2}w_1}{\mu_i-2d-\binom{n_2-1}{r-2}}\mathbf{1}_{n_2}^{\intercal}, \ldots , \frac{\binom{n_2-1}{r-2}w_{n_1}}{\mu_i-2d-\binom{n_2-1}{r-2} }\mathbf{1}_{n_2}^{\intercal} \right]^\intercal$. We claim that $(\mu_i, \mathbf{y}^{(i)})$ forms an eigenpair for $\mathcal{Q}(\mathcal{H})$.
		It is important to note that for any vertex $j \in [n_1]$,  $d_{\mathcal{H}}(j) = d_{\mathcal{H}_1}(j) + \binom{n_2}{r-1},$ and for any vertex $k \in [n_1(n_2+1)] \setminus [n_1],$ $d_{\mathcal{H}}(k)= d_{\mathcal{H}_2}(k) + \binom{n_2-1}{r-2}.$ 
		Proof of the claim is similar to the case of adjacency hypermatrix.
		
		\noindent Suppose $\eta$ is a non-main eigenvalue of $\mathcal{Q}(\mathcal{H}_2)$ with non-main eigenvector $\mathbf{z}$, then $\eta+\binom{n_2-1}{r-2}$ is also an eigenvalue of $\mathcal{Q}(\mathcal{H})$ with eigenvectors
		$\left[\begin{matrix}
			\mathbf{0} \\ \mathbf{z} \otimes \mathbf{e_j} 
		\end{matrix}\right], 1 \leq j \leq n_1.$
	\end{proof}

\end{document}
